% PAGES: 143-144
% Continues sec04_c3_1.tex (pages 141-142), which closed cleanly (Lemma
% 4.9's proof ends with the printed square at the bottom of page 142). No
% environment is open at the start of this file.
%
% Page 143 (folio 127): Theorem 4.8 (with proof), then a manually-typeset
% "Example:"/"Answer:" block (two nonsingular-matrix checks) -- same
% convention as sec04_c2_1.tex: the "Example:" run-in label is written out
% by hand (no closing mark after the statement), and \answerx is used for
% "Answer:" with a manual "\hfill$\square$" added just before
% \end{answerx}, matching the printed square at the bottom of page 143.
% This closes section 4.3.
%
% Page 144 (folio 128): section 4.4 "Numerical computation of eigenvalues"
% (complete, including footnote 5), then section 4.5 "Eigenvalues and
% eigenvectors of tridiagonal symmetric matrices" begins: matrix $B_N$
% (4.42), its entry definitions, and Lemma 4.10 opens with only its first
% sentence (the eigenvalue formula (4.43)) -- the page ends there. The
% source continues the SAME lemma statement with a second sentence (the
% eigenvector formula) on the next page (folio 129 / PDF page 145), still
% in italics, before "Proof." begins. The \begin{lemma}...\end{lemma}
% environment is therefore LEFT OPEN at the end of this file; it is closed
% at the top of sec04_c3_3.tex (pages 145-146), which completes the
% statement with (4.44) before opening "Proof."

\begin{theorem}
Any strictly column diagonally dominant matrix is nonsingular.
\end{theorem}

\begin{proof}
Let $A$ be a strictly column diagonally dominant matrix. From Lemma 4.9, it follows
that the matrix $A^t$ is strictly diagonally dominant. Thus, the matrix $A^t$ is
nonsingular, see Theorem 4.7, and therefore $0$ is not an eigenvalue of $A^t$. Since
$A$ and $A^t$ have the same eigenvalues, see Lemma 4.5, we conclude that $0$ is not an
eigenvalue of $A$, and therefore that $A$ is a nonsingular matrix.
\end{proof}

\par\medskip\noindent\textit{Example:}\ Show that the following matrices are
nonsingular:
\[
A_1 \;=\; \begin{pmatrix}
3 & 1 & -1 & 0\\
2 & -4 & 1 & 0.5\\
-2 & 1 & 5 & 0.5\\
0 & 2 & 0 & -2.5
\end{pmatrix};
\]
\[
A_2 \;=\; \begin{pmatrix}
-4 & 3 & -1 & 1\\
-1 & -5 & 1 & 4\\
1 & -1 & 3 & 0\\
0.5 & 0 & 0.5 & -6
\end{pmatrix}.
\]

\begin{answerx}
(i) The matrix $A_1$ is strictly diagonally dominant, since, for every row of $A_1$,
the absolute value of the main diagonal entry from the row is greater than the sum of
the absolute values of all the other entries in that row:
\begin{align*}
3 \;&>\; 1 + |{-1}| \;=\; 2;\\
|{-4}| = 4 \;&>\; 2 + 1 + 0.5 \;=\; 3.5;\\
5 \;&>\; |{-2}| + 1 + 0.5 \;=\; 3.5;\\
|{-2.5}| = 2.5 \;&>\; 2.
\end{align*}
Then, from Theorem 4.7, it follows that $A_1$ is a nonsingular matrix.

(ii) The matrix $A_2$ is strictly column diagonally dominant, since, for every column
of the matrix $A_2$, the absolute value of the main diagonal entry from the column is
strictly greater than the sum of the absolute values of all the other entries in that
column:
\begin{align*}
|{-4}| = 4 \;&>\; |{-1}| + 1 + 0.5 \;=\; 2.5;\\
|{-5}| = 5 \;&>\; 3 + |{-1}| \;=\; 4;\\
3 \;&>\; |{-1}| + 1 + 0.5 \;=\; 2.5;\\
|{-6}| = 6 \;&>\; 1 + 4 \;=\; 5.
\end{align*}

Then, from Theorem 4.8, it follows that $A_2$ is a nonsingular matrix.
Note that Theorem 4.7 cannot be used to show that the matrix $A_2$ is nonsingular,
since $A_2$ is not strictly diagonally dominant: for the first row of $A_2$, we find
that
\[
|{-4}| \;=\; 4 \;<\; 3 + |{-1}| + 1 \;=\; 5. \hfill\square
\]
\end{answerx}

\section{Numerical computation of eigenvalues}

Computing the eigenvalues of a matrix is a fundamental numerical linear algebra
problem, which further highlights the difference between numerical linear algebra and
classical linear algebra. For example, computing the roots of the characteristic
polynomial of a matrix to obtain its eigenvalues would be inefficient and prone to
large errors since polynomial root computation is an ill--conditioned problem.

The methods used in practice for computing eigenvalues are iterative methods that find
approximations of the eigenvalues of the matrix to desired tolerance. The simplest such
method is the Power Method which works best for symmetric positive definite matrices
whose largest eigenvalue has multiplicity 1, and the related Inverse Power Method and
Inverse Power Method with shifts. These methods are generally slowly convergent.

Efficient eigenvalue algorithms are variants of the QR Algorithm, which was named one
of the top ten algorithms of the 20--th century.\footnote{See Cipra \cite{10} at
\url{http://www.siam.org/pdf/news/637.pdf}} Based on the QR decomposition of a matrix
as the product of an orthogonal matrix Q and an upper triangular matrix R, the matrix
$A$ is iteratively transformed into $A_i = QR$ and then $A_{i+1} = RQ$, until
convergence to an upper triangular matrix (or to a diagonal matrix, if the matrix $A$
is symmetric) is achieved. More details on QR algorithms and their implementations can
be found, e.g., in Demmel \cite{13} and in Trefethen and Bau \cite{43}.

\section{Eigenvalues and eigenvectors of tridiagonal symmetric matrices}

Let
\begin{equation}
B_N \;=\; \begin{pmatrix}
2 & -1 & \dots & 0\\
-1 & \ddots & \ddots & \vdots\\
\vdots & \ddots & \ddots & -1\\
0 & \dots & -1 & 2
\end{pmatrix}
\end{equation}
be the $N \times N$ tridiagonal symmetric matrix whose only nonzero entries are
\begin{align*}
B_N(i,i) &= 2, \quad \forall\, i = 1:N;\\
B_N(i,i-1) &= -1, \quad \forall\, i = 2:N;\\
B_N(i,i+1) &= -1, \quad \forall\, i = 1:(N-1).
\end{align*}
The matrix $B_N$ appears often in practical applications, e.g., in the finite
difference solution of second order differential equations.

In this section, we find the eigenvalues and eigenvectors of the matrix $B_N$, and
then use them to compute the eigenvalues and eigenvectors of tridiagonal symmetric
matrices having the more general form (4.49). Note that, since the matrix $B_N$ is
symmetric, all its eigenvalues are real numbers; cf. Theorem 5.1.

\begin{lemma}
The matrix $B_N$ has the following $N$ different eigenvalues:
\begin{equation}
\mu_j \;=\; 2\left(1 - \cos\left(\frac{j\pi}{N+1}\right)\right), \quad \text{for}\ j = 1:N.
\end{equation}
% \begin{lemma} intentionally left open here (statement continues with the
% eigenvector formula on the next page); closed at the top of
% sec04_c3_3.tex.
